\documentclass[
paper=5.5in:8.5in,
%a5paper,
]{scrbook} % Document font size and paper size

\title{A Study in Scarlet}

\usepackage[namedChapters]{bindery}
\usepackage{dropcaps}
\usepackage{epistolary}
\usepackage{tabularx}

\setotherlanguage{french}
\setotherlanguage{latin}


\setdropcaps{Zallman Caps.ttf}


\DeclareTOCStyleEntry[linefill=\TOCLineLeaderFill,beforeskip=0.2cm]{chapter}{chapter}
\DeclareTOCStyleEntry[beforeskip=.3cm,dynnumwidth=true,entryformat=\Large\bfseries]{part}{part}

\setkomafont{chapter}{\chapterfont\Huge\bfseries}
\setkomafont{part}{\chapterfont\LARGE}
\setkomafont{partnumber}{\chapterfont\Huge\bfseries}

\hyphenation{}


\begin{document}
\frontmatter

  \begin{titlepage}
   \includepdf[width=\coverwidth]{titlepage2.jpg}
  \end{titlepage}



\tableofcontents



\mainmatter

\part[The Reminiscences of John H. Watson, M.D.]{Being a reprint from the Reminiscences of \textsc{John H. Watson, M.D.},\\Late of the Army Medical Department.}
\include{chapters/01.tex}
\include{chapters/02.tex}
\include{chapters/03.tex}
\include{chapters/04.tex}
\include{chapters/05.tex}
\include{chapters/06.tex}
\include{chapters/07.tex}
\part{The Country of the Saints}
\include{chapters/08.tex}
\include{chapters/09.tex}
\include{chapters/10.tex}
\include{chapters/11.tex}
\include{chapters/12.tex}
\include{chapters/13.tex}
\include{chapters/14.tex}


\KOMAoptions{headings=openleft}
\chapter*{Colophon}

\centering
\vfill
\begin{minipage}{\textwidth}
After first appearing in the 1887 edition of \textit{Beeton's Christmas Annual}, \textit{A Study in Scarlet} was published as a standalone novel in July 1888 by Ward, Lock \& Co. in London (UK). It was Arthur Conan Doyle's (1859--1930) first novel, and introduced Sherlock Holmes and John Watson.

Title page illustration by David Henry Friston (1820--1906) is from the original publication in \textit{Beeton's Christmas Annual}.
\end{minipage}
\vfill
\rule{0.5\textwidth}{.4pt}
\vfill
\begin{minipage}{\textwidth}
EB Garamond is Georg Mayr-Duffner's free and open source implementation of Claude Garamont’s famous humanist typefaces from the mid-sixteenth century. This digital version reproduces the original design by Claude Garamont closely: the source for the letterforms is a scan of a specimen known as the <Berner specimen>, which was composed in 1592 by Conrad Berner, the son-in-law of Christian Egenolff and his successor at the Egenolff print office. Dropcaps are set in Dieter Steffman's <Zallman Caps>. Title page (author credit only) is set in <Bolton,> by Paul Lloyd.
\end{minipage}
github.com/georgd/EB-Garamond\\
steffmann.de\\
moorstation.org/typoasis/designers/lloyd
\vfill
\rule{0.5\textwidth}{.4pt}
\vfill
\begin{minipage}{\textwidth}
This typeset is dedicated to the public domain under a Creative Commons CC0 1.0 Universal deed: creativecommons.org/publicdomain/zero/1.0/
\end{minipage}
\vfill
\rule{0.5\textwidth}{.4pt}
\vfill
Typeset in \LaTeX{}. Last revised \today.
\thispagestyle{empty}
\end{document}